%% 01_problem_formulation.tex -- Problem Formulation and Constraint Specifications
%% -----------------------------------------------------------------------------

\chapter{Problem Formulation}

\section{Introduction and Academic Context}
Class scheduling (also known as university course timetabling) represents one of the most widely studied and computationally challenging combinatorial optimization problems in artificial intelligence and operations research. The objective is to allocate a set of lectures, classrooms, timeslots, and instructors into a conflict-free, highly satisfying weekly timetable subject to a rigorous set of institutional constraints.

In this project, we address a large-scale university timetabling problem with the structural specifications summarized in Table~\ref{tab:problem_dimensions}.

\begin{table}[H]
  \centering
  \small
  \caption{Problem Dimensions and Parameter Specifications}
  \label{tab:problem_dimensions}
  \begin{tabularx}{\linewidth}{L{4.2cm} L{3.0cm} X}
    \toprule
    \rowcolor{PrismHeader}
    \color{PrismWhite}\textbf{Parameter} & 
    \color{PrismWhite}\textbf{Notation / Value} & 
    \color{PrismWhite}\textbf{Description / Domain} \\
    \midrule
    Courses & $N_c = 60$ & Distinct courses $\{c_{01}, c_{02}, \dots, c_{60}\}$. \\
    Sections & $N_{sec} = 480$ & 8 sections per course ($60 \times 8 = 480$). \\
    Sessions (Lectures) & $N_s = 960$ & 2 sessions per section ($480 \times 2 = 960$ timeslots total). \\
    Professors & $N_p = 60$ & Instructors $\{p_{01}, p_{02}, \dots, p_{60}\}$, max 20 sessions/week. \\
    Classrooms & $N_r = 35$ & Available physical rooms $\{rm_{01}, \dots, rm_{35}\}$. \\
    Room Sizes & $s_r \in \{0, 1, 2\}$ & 0: Small, 1: Average, 2: Large. \\
    Registration Sizes & $s_{sec} \in \{0, 1, 2\}$ & Enrollment capacity requirement (0: Small, 1: Avg, 2: Large). \\
    Teaching Days & $D = 5$ & Days of the week: Monday (0) to Friday (4). \\
    Timeslots per Day & $T = 6$ & 6 slots per day (Slots 1--3: Morning, Slots 4--6: Afternoon). \\
    Total Available Slots & $D \times T \times N_r = 1050$ & Room utilization is high ($\approx 960 / 1050 = 91.43\%$). \\
    \bottomrule
  \end{tabularx}
\end{table}

\section{Mathematical Modeling}
Let $\mathcal{S} = \{s_1, s_2, \dots, s_{960}\}$ be the set of all sessions. Each session $s_i$ belongs to a specific section $sec(s_i)$ and course $c(s_i)$, with an enrollment size $s_{sec}(s_i) \in \{0, 1, 2\}$.
A complete schedule solution is a mapping:
\begin{equation}
  \sigma: \mathcal{S} \rightarrow \mathcal{D} \times \mathcal{T} \times \mathcal{R} \times \mathcal{P}
\end{equation}
where for each session $s_i$, the schedule assigns:
\begin{itemize}[leftmargin=1.5em]
  \item $d(s_i) \in \mathcal{D} = \{0, 1, 2, 3, 4\}$ (Monday to Friday),
  \item $t(s_i) \in \mathcal{T} = \{0, 1, 2, 3, 4, 5\}$ (Periods 1 to 6),
  \item $r(s_i) \in \mathcal{R} = \{0, 1, \dots, 34\}$ (Rooms $rm_{01}$ to $rm_{35}$),
  \item $p(s_i) \in \mathcal{P} = \{0, 1, \dots, 59\}$ (Professors $p_{01}$ to $p_{60}$).
\end{itemize}

\section{Constraint Hierarchy}
Constraints are categorized into \textbf{Hard Constraints} (mandatory requirements that define feasibility), \textbf{Soft Constraints} (quality metrics to be optimized), and \textbf{Bonus Objectives} (preference satisfaction).

\subsection{Hard Constraints (Feasibility Boundary)}
Any valid academic schedule must satisfy the following six hard constraints without violation:
\begin{enumerate}[leftmargin=2.0em]
  \item \textbf{H1 --- Room/Section Capacity Compatibility:}
  A classroom must be large enough to accommodate the registered section size:
  \begin{equation}
    \text{Size}(r(s_i)) \ge s_{sec}(s_i), \quad \forall s_i \in \mathcal{S}
  \end{equation}
  Rooms of size 2 can host sizes $\{0, 1, 2\}$; rooms of size 1 can host $\{0, 1\}$; rooms of size 0 can only host size $0$.

  \item \textbf{H2 --- Section Day Separation (Non-Adjacent Days):}
  Each section has exactly 2 sessions per week. They \textbf{cannot} occur on the same day or on adjacent days:
  \begin{equation}
    d(s_{k, 1}) \ne d(s_{k, 2}) \quad \text{and} \quad |d(s_{k, 1}) - d(s_{k, 2})| > 1, \quad \forall sec_k
  \end{equation}
  Valid pairs: (Mon, Wed), (Mon, Thu), (Mon, Fri), (Tue, Thu), (Tue, Fri), (Wed, Fri).

  \item \textbf{H3 --- Professor Conflict \& Workload Ceiling:}
  A professor cannot teach two sessions simultaneously:
  \begin{equation}
    \sum_{s_i \in \mathcal{S}} \mathbf{1}_{\{p(s_i) = p \land d(s_i) = d \land t(s_i) = t\}} \le 1, \quad \forall p \in \mathcal{P}, \forall d \in \mathcal{D}, \forall t \in \mathcal{T}
  \end{equation}
  Furthermore, no professor may exceed 20 teaching sessions per week:
  \begin{equation}
    \sum_{s_i \in \mathcal{S}} \mathbf{1}_{\{p(s_i) = p\}} \le 20, \quad \forall p \in \mathcal{P}
  \end{equation}

  \item \textbf{H4 --- Classroom Availability:}
  All sessions must be placed in valid, operational classrooms ($r(s_i) \in \mathcal{R}$).

  \item \textbf{H5 --- Section Room Uniqueness:}
  A section cannot be scheduled in two different rooms at the exact same timeslot.

  \item \textbf{H6 --- Classroom Overlap Collision:}
  A classroom cannot host more than one session during any single timeslot:
  \begin{equation}
    \sum_{s_i \in \mathcal{S}} \mathbf{1}_{\{r(s_i) = r \land d(s_i) = d \land t(s_i) = t\}} \le 1, \quad \forall r \in \mathcal{R}, \forall d \in \mathcal{D}, \forall t \in \mathcal{T}
  \end{equation}
\end{enumerate}

\subsection{Soft Constraints (Timetable Desirability)}
\begin{enumerate}[leftmargin=2.0em]
  \item \textbf{S1 --- Professor Schedule Compactness (No Idle Gaps):}
  Professor daily schedules should not scatter with vacant intermediate periods (e.g., slot 1 and slot 6 on the same day).
  \item \textbf{S2 --- Professor Half-Day Consistency:}
  Instructors should teach exclusively in the morning ($t \in \{0, 1, 2\}$) or afternoon ($t \in \{3, 4, 5\}$) on any given day.
  \item \textbf{S3 --- Course Half-Day Consistency:}
  All sections of a course should ideally be held either exclusively in the morning or exclusively in the afternoon.
  \item \textbf{S4 --- Room Capacity Efficiency:}
  Minimize wasted space ($\text{Size}(r(s_i)) - s_{sec}(s_i)$) so that large rooms are not squandered on small sections.
\end{enumerate}

\subsection{Bonus Constraints}
\begin{enumerate}[leftmargin=2.0em]
  \item \textbf{B1 --- Faculty Preference Satisfaction:}
  Each professor lists their top 3 favorite courses. The algorithm optimizes so that every professor is assigned to teach at least one favorite course.
  \item \textbf{B2 --- Instructor Qualification Domain:}
  A professor should not be assigned to courses outside their verified expertise list.
\end{enumerate}
